\documentclass[11pt,a4paper]{article}

\usepackage[margin=1in]{geometry}
\usepackage{amsmath,amssymb,amsthm,mathtools}
\usepackage{comment}
\usepackage{enumitem}
\usepackage[hidelinks]{hyperref}
\usepackage{xurl}
\urlstyle{tt}
\renewcommand\UrlBreaks{\do\.\do\@\do\\\do\/\do\!\do\_\do\?\do\&\do\=\do\-\do\:\do\;}

\newtheorem{theorem}{Theorem}[section]
\newtheorem{proposition}[theorem]{Proposition}
\newtheorem{lemma}[theorem]{Lemma}
\newtheorem{corollary}[theorem]{Corollary}
\theoremstyle{definition}
\newtheorem{definition}[theorem]{Definition}
\newtheorem{example}[theorem]{Example}
\newtheorem{remark}[theorem]{Remark}
\newtheorem{assumption}[theorem]{Assumption}

\newcommand{\R}{\mathbb{R}}
\newcommand{\C}{\mathbb{C}}
\newcommand{\N}{\mathbb{N}}
\newcommand{\Z}{\mathbb{Z}}
\newcommand{\K}{\mathbb{K}}
\newcommand{\HH}{\mathbb{H}}
\newcommand{\Tpoly}{T_{\mathrm{poly}}}
\newcommand{\Dpoly}{D_{\mathrm{poly}}}
\newcommand{\Cinfty}{C^{\infty}}
\newcommand{\Sch}[2]{[\,#1,#2\,]_{\mathrm{SN}}}
\newcommand{\Ger}[2]{[\,#1,#2\,]_{\mathrm{G}}}
\newcommand{\defi}[1]{\emph{#1}}

\title{Kontsevich Deformation Quantization in Lean\\
{\large Problem statement and a formalization path}}
\author{Track 4 working note}
\date{\today}

\begin{document}
\maketitle

\begin{abstract}
This note states the deformation quantization problem solved by
Kontsevich's formality theorem, fixes notation for every structure
involved, and lays out a layered path for formalizing the result in
Lean 4 / Mathlib. The note is intentionally minimal. It records what
must be defined, what is already available in Mathlib, what is missing,
and which subproblems can be attacked first. References are cited by
paper, section, and result number throughout.
\end{abstract}




\section{Background and goal}\label{sec:background}

Let $M$ be a smooth manifold. Write $T_M$ for its tangent bundle:
\[
T_M := \bigsqcup_{p \in M} T_pM,
\]
where $T_pM$ is the tangent vector space at $p$. A smooth vector field
is a section of $T_M$, i.e.\ an element of $\Gamma(T_M)$. A bivector field is a section of $\Lambda^2 T_M$, i.e.\ an element of
$\Gamma(\Lambda^2 T_M)$.

A \defi{Poisson manifold} is a pair $(M,\pi)$ where $M$ is a smooth
manifold and $\pi \in \Gamma(\Lambda^2 T_M)$ is a bivector field whose
Schouten--Nijenhuis self-bracket vanishes:
\[
\Sch{\pi}{\pi}=0.
\]
This condition is equivalent to the Jacobi identity for the bracket
\[
\{f,g\} := \pi(df, dg), \qquad f,g \in \Cinfty(M).
\]
Thus $\{\,\cdot\,,\,\cdot\,\}$ is a Poisson bracket on $\Cinfty(M)$.
In local coordinates, if
\[
\pi=\frac12\sum_{i,j}\pi^{ij}(x)\,\partial_i\wedge\partial_j,
\]
then $\Sch{\pi}{\pi}=0$ is equivalent to
\[
\sum_{\ell}
\left(
\pi^{i\ell}\partial_\ell\pi^{jk}
+ \pi^{j\ell}\partial_\ell\pi^{ki}
+ \pi^{k\ell}\partial_\ell\pi^{ij}
\right)=0
\qquad \text{for all } i,j,k.
\]

The problem posed by physics is to associate, to each
$(M,\pi)$, an associative noncommutative product $\star$ on
$\Cinfty(M)[[\hbar]]$ whose semiclassical limit is $\pi$. The precise
notion is a star product (Bayen--Flato--Fronsdal--Lichnerowicz--Sternheimer
\cite{BFFLS}).

\begin{definition}\label{def:star}
A \defi{star product} on $(M,\pi)$ is an $\R[[\hbar]]$-bilinear map
\[
\star \colon \Cinfty(M)[[\hbar]] \otimes_{\R[[\hbar]]} \Cinfty(M)[[\hbar]]
\longrightarrow \Cinfty(M)[[\hbar]],
\quad
f \star g = fg + \sum_{k \ge 1} \hbar^{k}\, B_{k}(f,g),
\]
such that:
\begin{enumerate}[label=(\roman*),leftmargin=2em]
  \item each $B_{k}$ is a bidifferential operator on $M$;
  \item $\star$ is associative;
  \item $B_{1}(f,g) - B_{1}(g,f) = \{f,g\}$.
\end{enumerate}
Two star products $\star, \star'$ are \defi{gauge equivalent} if there
exists $T = \mathrm{id} + \sum_{k\ge 1}\hbar^{k} T_{k}$ with $T_{k}$
differential operators and $T(f \star g) = T f \star' T g$.
\end{definition}

\begin{example}[Standard Poisson structure]
\label{ex:standard-poisson}
Let $\R^{2m}$ have coordinates
$(q_1,\dots,q_m,p_1,\dots,p_m)$. The bivector field
\[
\pi_0 := \sum_{i=1}^{m}
\frac{\partial}{\partial q_i}\wedge
\frac{\partial}{\partial p_i}
\]
is Poisson. Its coefficients are constant, hence
$\Sch{\pi_0}{\pi_0}=0$. The induced bracket is
\[
\{f,g\}
=\sum_{i=1}^{m}
\left(
\frac{\partial f}{\partial q_i}\frac{\partial g}{\partial p_i}
-\frac{\partial f}{\partial p_i}\frac{\partial g}{\partial q_i}
\right).
\]
\end{example}

\begin{example}[The Moyal product]
\label{ex:moyal-background}
Fix an antisymmetric constant matrix $\pi^{ij}=-\pi^{ji}$ on $\R^d$
and set
$\pi = \tfrac{1}{2}\sum_{i,j}\pi^{ij}\,\partial_i\wedge\partial_j$.
The Moyal product is
\[
f\star g
=m\left(
\exp\left(
\frac{\hbar}{2}\sum_{i,j}\pi^{ij}\partial_i\otimes\partial_j
\right)(f\otimes g)
\right),
\]
where $m(f\otimes g)=fg$. Its first terms are
\[
f\star g
=fg+\frac{\hbar}{2}\sum_{i,j}\pi^{ij}(\partial_i f)(\partial_j g)
+O(\hbar^2),
\]
so
\[
f\star g-g\star f=\hbar\{f,g\}+O(\hbar^2).
\]
\end{example}

The goal of Kontsevich's formality theorem is to show that such a
$\star$ exists for every Poisson structure, and to classify them up to
gauge equivalence by Poisson structures up to formal diffeomorphism.

\section{The two graded Lie algebras}\label{sec:two-gla}

Kontsevich's framework expresses the problem as a comparison of two
$L_\infty$-algebras built from $M$.

\subsection{Polyvector fields}\label{ssec:Tpoly}

\begin{definition}\label{def:Tpoly}
For a smooth manifold $M$, set
\[
\Tpoly^{n}(M) := \Gamma\bigl(\Lambda^{n+1} T_M\bigr), \qquad n \ge -1.
\]
Equipped with the Schouten--Nijenhuis bracket $\Sch{\cdot}{\cdot}$ and
zero differential, $\Tpoly(M)$ is a graded Lie algebra concentrated in
degrees $\ge -1$.
\end{definition}

A Maurer--Cartan element of $\Tpoly(M)[[\hbar]]$ in degree~$1$ is an
element $\pi = \hbar \pi_{1} + \hbar^{2} \pi_{2} + \cdots$ with
$\pi_{k} \in \Gamma(\Lambda^{2} T_M)$ and $\Sch{\pi}{\pi} = 0$.

\begin{example}[Polyvector fields on $\R^2$]
\label{ex:polyvectors-r2}
For $M=\R^2$ with coordinates $(x,y)$,
\[
\Tpoly^{-1}(M)=\Cinfty(M),\qquad
\Tpoly^0(M)=\Gamma(T_M),\qquad
\Tpoly^1(M)=\Gamma(\Lambda^2 T_M).
\]
A vector field is, for example,
\[
X=x\partial_x+y\partial_y.
\]
A bivector field is, for example,
\[
\pi=f(x,y)\partial_x\wedge\partial_y.
\]
Every bivector field on a surface is Poisson, because
$\Sch{\pi}{\pi}$ is a trivector field and $\Lambda^3 T_M=0$.
\end{example}

\begin{example}[Polyvector fields on the Riemann sphere]
\label{ex:polyvectors-cp1}
Let $M = S^{2} \cong \mathbb{CP}^{1}$ viewed as a real $2$-manifold.
Cover $M$ by two stereographic charts:
\[
U_{N} = M \setminus \{S\}, \ \ z = x+iy \in \C;
\qquad
U_{S} = M \setminus \{N\}, \ \ w = u+iv \in \C;
\]
with transition $w = 1/z$ on $U_{N} \cap U_{S}$.

The tangent bundle $T_M$ is a real rank-$2$ vector bundle, locally
trivialised by $(\partial_x, \partial_y)$ on $U_{N}$ and by
$(\partial_u, \partial_v)$ on $U_{S}$. As a bundle over $S^{2}$ it
is nontrivial: every smooth global vector field on $S^{2}$ has a
zero (Hairy Ball Theorem). The polyvector-field grading is
\[
\Tpoly^{-1}(M) = \Cinfty(M),\qquad
\Tpoly^{0}(M) = \Gamma(T_M),\qquad
\Tpoly^{1}(M) = \Gamma(\Lambda^{2} T_M),
\]
and $\Tpoly^{n}(M) = 0$ for $n \ge 2$ because
$\Lambda^{k} T_{p}M = 0$ for $k > \dim_{\R} M = 2$.

The bundle $\Lambda^{2}T_M$ has real rank $\binom{2}{2} = 1$; a global
section is a bivector field. On the two charts,
\[
\pi|_{U_{N}} = f(x,y)\, \partial_x \wedge \partial_y,
\qquad
\pi|_{U_{S}} = g(u,v)\, \partial_u \wedge \partial_v,
\]
with $f(x,y) = \det J(x,y) \cdot g(u(x,y), v(x,y))$ on the overlap,
where $J = \partial(u,v)/\partial(x,y)$ is the Jacobian of the change
of coordinates.

A canonical choice of $\pi \in \Gamma(\Lambda^{2}T_M)$ comes from the
Fubini--Study K\"ahler form
\[
\omega_{\mathrm{FS}}
= \frac{i}{2}\,\frac{dz \wedge d\bar z}{(1+|z|^{2})^{2}}
= \frac{dx \wedge dy}{(1+x^{2}+y^{2})^{2}},
\]
inverted to the bivector
\[
\pi_{\mathrm{FS}}|_{U_{N}}
= (1+x^{2}+y^{2})^{2}\, \partial_x \wedge \partial_y,
\qquad
\pi_{\mathrm{FS}}|_{U_{S}}
= (1+u^{2}+v^{2})^{2}\, \partial_u \wedge \partial_v.
\]
Both local expressions come from inverting the same intrinsic
$\omega_{\mathrm{FS}}$, so they glue on $U_{N} \cap U_{S}$.

Poisson-ness is automatic on $S^{2}$: for any
$\pi \in \Gamma(\Lambda^{2}T_M)$, $\Sch{\pi}{\pi} \in
\Gamma(\Lambda^{3} T_M) = 0$ by the same dimension argument. So
$(S^{2}, \pi_{\mathrm{FS}})$ is a Poisson manifold, with bracket
\[
\{f, g\}(z)
= (1+|z|^{2})^{2}
\bigl(\partial_x f \cdot \partial_y g
    - \partial_y f \cdot \partial_x g\bigr)
\]
on $U_{N}$. This is the Poisson bracket of the Fubini--Study
symplectic structure and the input data for a Kontsevich star
product on $S^{2}$.

The holomorphic tangent bundle $T_{\mathbb{CP}^{1}}^{1,0} \cong
\mathcal{O}(2)$ has complex rank $1$, so
$\Lambda^{2} T_{\mathbb{CP}^{1}}^{1,0} = 0$. There is no non-trivial
holomorphic Poisson structure on $\mathbb{CP}^{1}$. The Kontsevich
framework as stated in Section~\ref{sec:background} uses the real
tangent bundle for this reason.
\end{example}

\subsection{Polydifferential operators}\label{ssec:Dpoly}

\begin{definition}\label{def:Dpoly}
Let $A := \Cinfty(M)$. Define
\[
\Dpoly^{n}(M) := \bigl\{\,D \colon A^{\otimes (n+1)} \to A \;\big|\;
D \text{ is polydifferential}\,\bigr\}, \qquad n \ge -1.
\]
The Hochschild differential $d_{\mathrm{H}} = [\mu, \,\cdot\,]_{\mathrm{G}}$,
where $\mu$ is the pointwise product, and the Gerstenhaber bracket
$\Ger{\cdot}{\cdot}$ make $\Dpoly(M)$ into a differential graded Lie
algebra (Gerstenhaber \cite[\S6]{Gerstenhaber1963}).
\end{definition}

A Maurer--Cartan element of $\hbar\Dpoly(M)[[\hbar]]$ in degree~$1$ is
$\pi = \hbar B_{1} + \hbar^{2} B_{2} + \cdots$ with each $B_{k}$
bidifferential and
\[
d_{\mathrm{H}}\pi + \tfrac{1}{2}\Ger{\pi}{\pi}=0.
\]
This equation is equivalent to associativity of the deformed product
$B := \mu + \pi = \mu + \hbar B_{1} + \hbar^{2} B_{2} + \cdots$. So
MC elements of $\Dpoly$ are exactly star products.

\begin{example}[A polydifferential operator]
\label{ex:polydifferential}
Let $A=\Cinfty(\R)$. The map
\[
D(f,g)=f'g''
\]
is a bidifferential operator, hence an element of $\Dpoly^1(\R)$.
The ordinary product $\mu(f,g)=fg$ is also an element of $\Dpoly^1(M)$.
The Gerstenhaber square measures the associator:
\[
\tfrac{1}{2}\,\Ger{\mu}{\mu}(f,g,h) = (fg)h - f(gh),
\]
so $\Ger{\mu}{\mu}=0$ is equivalent to associativity.
\end{example}

\subsection{The HKR map}\label{ssec:HKR}

\begin{theorem}[Hochschild--Kostant--Rosenberg]\label{thm:HKR}
The map $U_{\mathrm{HKR}} \colon \Tpoly(M) \to \Dpoly(M)$ defined by
\[
U_{\mathrm{HKR}}(\xi_{1} \wedge \cdots \wedge \xi_{n})(f_{1}, \dots, f_{n})
:= \frac{1}{n!} \sum_{\sigma \in S_{n}} \mathrm{sgn}(\sigma)\,
\xi_{\sigma(1)}(f_{1}) \cdots \xi_{\sigma(n)}(f_{n})
\]
is a quasi-isomorphism of cochain complexes.
\end{theorem}

The map $U_{\mathrm{HKR}}$ is \emph{not} a morphism of dg Lie algebras.
The failure is the source of the problem. Kontsevich's theorem upgrades
$U_{\mathrm{HKR}}$ to a structure that does respect the brackets, in a
relaxed sense.

\begin{example}[HKR on a decomposable bivector]
\label{ex:hkr}
For vector fields $X,Y$,
\[
U_{\mathrm{HKR}}(X\wedge Y)(f,g)
=\frac12\left(X(f)Y(g)-Y(f)X(g)\right).
\]
Thus a bivector field becomes an antisymmetric bidifferential operator.
\end{example}

\section{\texorpdfstring{$L_\infty$}{L-infinity}-algebras and the formality theorem}\label{sec:Linfty}

\begin{definition}\label{def:Linfty}
An \defi{$L_\infty$-algebra} on a graded $\K$-module $L$ is a family of
graded antisymmetric maps
\[
Q_{n} \colon L^{\wedge n} \to L, \qquad |Q_{n}| = 2 - n,\quad n \ge 1,
\]
satisfying the generalized Jacobi identities, for every $n \ge 1$,
\begin{equation*}
\sum_{i+j=n+1}\;\sum_{\sigma \in \mathrm{Sh}(i,n-i)}
\mathrm{sgn}(\sigma)\, \varepsilon(\sigma)\,
Q_{j}\!\bigl(Q_{i}(x_{\sigma(1)}, \dots, x_{\sigma(i)}),\;
x_{\sigma(i+1)}, \dots, x_{\sigma(n)}\bigr) \;=\; 0,
\end{equation*}
where $\varepsilon(\sigma)$ is the Koszul sign and $\mathrm{Sh}(i, n-i)$
is the set of $(i, n-i)$-shuffles.
\end{definition}

A differential graded Lie algebra $(L, d, [\cdot,\cdot])$ is the special
case $Q_{1} = d$, $Q_{2} = [\cdot,\cdot]$, $Q_{n \ge 3} = 0$.

\begin{example}[A dg Lie algebra as an $L_\infty$-algebra]
\label{ex:dgla-linfty}
If $(L,d,[\,\cdot\,,\,\cdot\,])$ is a differential graded Lie algebra,
then it becomes an $L_\infty$-algebra by setting
\[
Q_1=d,\qquad Q_2=[\,\cdot\,,\,\cdot\,],\qquad Q_n=0\quad(n\ge3).
\]
The $n=1$ identity says $d^2=0$. The $n=2$ identity says that $d$ is
compatible with the bracket. The $n=3$ identity says the graded Jacobi
identity.
\end{example}

\begin{definition}\label{def:Linfty-morphism}
An \defi{$L_\infty$-morphism} $U \colon (L, Q) \to (L', Q')$ is a family
of graded antisymmetric maps
\[
U_{n} \colon L^{\wedge n} \to L', \qquad |U_{n}| = 1 - n,\quad n \ge 1,
\]
satisfying the compatibility identities of \cite[\S3.4]{Kontsevich2003}.
The first component $U_{1}$ is a chain map $(L, Q_{1}) \to (L', Q'_{1})$.
A morphism $U$ is a \defi{quasi-isomorphism} if $U_{1}$ is.
\end{definition}

\begin{definition}\label{def:MC}
A \defi{Maurer--Cartan element} of an $L_\infty$-algebra $(L, Q)$ over
$\hbar \K[[\hbar]]$ is $\alpha \in L^{1} \otimes \hbar \K[[\hbar]]$ such
that
\[
\sum_{n \ge 1} \frac{1}{n!}\, Q_{n}(\alpha, \dots, \alpha) = 0.
\]
\end{definition}

\begin{example}[Maurer--Cartan in a dg Lie algebra]
\label{ex:mc-dgla}
For a dg Lie algebra, the Maurer--Cartan equation is
\[
d\alpha+\frac12[\alpha,\alpha]=0.
\]
For $\Tpoly(M)$ the differential is zero, so the equation becomes
$\Sch{\pi}{\pi}=0$. Thus Poisson bivectors are exactly
Maurer--Cartan elements of $\Tpoly(M)$. For $\Dpoly(M)$, the same
equation is associativity of the deformed product.
\end{example}

An $L_\infty$-morphism $U$ transports MC elements: if $\alpha$ is MC in
$L$, then $\sum_{n \ge 1} \tfrac{1}{n!} U_{n}(\alpha, \dots, \alpha)$
is MC in $L'$.

\begin{theorem}[Kontsevich formality {\cite[Thm.~6.4]{Kontsevich2003}}]
\label{thm:formality}
For every smooth manifold $M$ there exists an $L_\infty$-quasi-isomorphism
\[
U \colon \Tpoly(M) \longrightarrow \Dpoly(M)
\]
whose first Taylor component $U_{1}$ equals $U_{\mathrm{HKR}}$ of
Theorem~\ref{thm:HKR}.
\end{theorem}

\begin{corollary}[Quantization]\label{cor:quantization}
For every Poisson structure $\pi$ on $M$, the formal series
\[
f \star g := fg + \sum_{n \ge 1} \frac{\hbar^{n}}{n!}\,
U_{n}(\pi, \dots, \pi)(f, g)
\]
defines a star product on $(M, \pi)$, and the assignment
$\pi \mapsto \star$ induces a bijection between formal Poisson
structures modulo formal diffeomorphism and star products modulo gauge
equivalence.
\end{corollary}

\begin{example}[How formality produces the star product]
\label{ex:formality-star}
For a Poisson bivector $\pi$, Kontsevich's maps $U_n$ produce
\[
f\star g
=fg+\sum_{n\ge1}\frac{\hbar^n}{n!}U_n(\pi,\dots,\pi)(f,g).
\]
The $L_\infty$ identities say that the condition
$\Sch{\pi}{\pi}=0$ is transported to the associativity of $\star$.
\end{example}

\section{The graphical formula on \texorpdfstring{$\R^{d}$}{R\^d}}
\label{sec:graphs}

For $M = \R^{d}$ Kontsevich constructs $U$ by an explicit sum over
graphs. This is the formula that has to be implemented.

\begin{definition}[Kontsevich admissible graph {\cite[\S6.1]{Kontsevich2003}}]
\label{def:graph}
For $n, m \ge 0$, an \defi{admissible graph} $\Gamma$ of type $(n, m)$
is a finite directed graph with:
\begin{itemize}[leftmargin=2em]
  \item vertex set $V(\Gamma) = V_{1} \sqcup V_{2}$, where
        $V_{1} = \{1, \dots, n\}$ are \defi{aerial} vertices and
        $V_{2} = \{\bar 1, \dots, \bar m\}$ are \defi{ground} vertices;
  \item each aerial vertex has exactly two outgoing edges, labelled
        $(\cdot, 1)$ and $(\cdot, 2)$;
  \item ground vertices have no outgoing edges;
  \item no edge ends at its starting vertex; no two edges share both
        endpoints.
\end{itemize}
Write $G_{n,m}$ for the (finite) set of admissible graphs of type
$(n, m)$. Each $\Gamma \in G_{n,m}$ has $2n$ edges.
\end{definition}

\begin{remark}\label{rem:bivectors-only}
The definition above matches Kontsevich's bivector case: every aerial
vertex has outdegree $2$, so $B_{\Gamma}$ takes $n$ bivector fields as
inputs. The general formality theorem allows aerial vertices of
arbitrary outdegree $a_{k}+1$, with input a polyvector in
$\Gamma(\Lambda^{a_{k}+1} T_M)$; in that case
$|E(\Gamma)|=\sum_{k}(a_{k}+1)$. We restrict to outdegree $2$
throughout because it suffices to construct the star product on a
Poisson manifold.
\end{remark}

\begin{example}[The first admissible graph]
\label{ex:first-graph}
Take $n=1$ and $m=2$. There is one aerial vertex $1$ and two ground
vertices $\bar1,\bar2$. Let the two outgoing edges from $1$ end at
$\bar1$ and $\bar2$. This graph is the first-order graph in the
Kontsevich formula.
\end{example}

\begin{definition}[Polydifferential operator $B_\Gamma$]
\label{def:BGamma}
For $\Gamma \in G_{n,m}$, bivector fields $\gamma_{1}, \dots, \gamma_{n}$
on $\R^{d}$, and functions $f_{1}, \dots, f_{m} \in \Cinfty(\R^{d})$, set
\[
B_{\Gamma}(\gamma_{1}, \dots, \gamma_{n})(f_{1}, \dots, f_{m})
:= \sum_{I \colon E(\Gamma) \to \{1, \dots, d\}}
\prod_{k=1}^{n} \Bigl(\prod_{e \in \mathrm{in}(k)} \partial_{I(e)}\Bigr)
\gamma_{k}^{I(e_{k,1}) I(e_{k,2})}
\cdot
\prod_{j=1}^{m} \Bigl(\prod_{e \in \mathrm{in}(\bar j)} \partial_{I(e)}\Bigr) f_{j},
\]
where $\mathrm{in}(v)$ is the multiset of edges ending at $v$ and
$(e_{k,1}, e_{k,2})$ are the two outgoing edges at aerial vertex $k$.
\end{definition}

\begin{example}[The operator of the first graph]
\label{ex:first-operator}
For the graph in Example~\ref{ex:first-graph} and a bivector
$\pi = \tfrac{1}{2}\sum_{i,j}\pi^{ij}\,\partial_i\wedge\partial_j$
with $\pi^{ij}=-\pi^{ji}$,
\[
B_\Gamma(\pi)(f,g)
=\sum_{i,j}\pi^{ij}(\partial_i f)(\partial_j g).
\]
The two labelings of the outgoing edges $(1\to\bar 1,\,1\to\bar 2)$
versus $(1\to\bar 2,\,1\to\bar 1)$ contribute $\pi^{ij}$ and
$\pi^{ji}$, and the antisymmetry of $\pi^{ij}$ folds them into a
single term. This is the first-order bidifferential operator attached
to the Poisson bracket.
\end{example}

\begin{definition}[Configuration space and weight]\label{def:weight}
Let $\HH := \{z \in \C : \Im z > 0\}$. For $2n + m - 2 \ge 0$, set
\[
C_{n,m} := \bigl\{ (p_{1}, \dots, p_{n}; q_{1}, \dots, q_{m}) \in
\HH^{n} \times \R^{m} \;\big|\; \text{all distinct},\ q_{1} < \cdots < q_{m} \bigr\}
\;\bigm/\; G_{\mathrm{aff}},
\]
where $G_{\mathrm{aff}} = \{ z \mapsto az + b : a > 0,\ b \in \R\}$.
Let $\bar C_{n,m}$ denote its Fulton--MacPherson--Axelrod--Singer
compactification \cite[\S5]{Kontsevich2003}, a manifold with corners of
dimension $2n + m - 2$.

For an edge $e = (v_{1}, v_{2})$ of an admissible graph and a
configuration in $C_{n,m}$, let
$\phi_{e} \in \R / 2\pi\Z$ be the hyperbolic angle from the geodesic
ray $v_{1} \to \infty$ in the Poincar\'e model of $\HH$ to the
geodesic from $v_{1}$ to $v_{2}$.

For $\Gamma \in G_{n,m}$ with $|E(\Gamma)| = 2n + m - 2$ define
\[
w_{\Gamma} :=
\frac{1}{(2\pi)^{2n+m-2}}
\int_{\bar C_{n,m}^{+}} \bigwedge_{e \in E(\Gamma)} d\phi_{e}\,,
\]
where $\bar C_{n,m}^{+}$ is the connected component in which the
$q_{j}$ are ordered $q_{1}<\cdots<q_{m}$ on $\R$ and the wedge is
taken in a fixed total order on $E(\Gamma)$. For all other $\Gamma$
set $w_{\Gamma} := 0$. Since our aerial vertices have outdegree $2$,
we have $|E(\Gamma)| = 2n$, and the condition
$|E(\Gamma)|=2n+m-2$ forces $m = 2$: nonzero weights occur only for
$\Gamma \in G_{n,2}$.
\end{definition}

\begin{example}[The configuration space $C_{1,2}$]
\label{ex:config-12}
For $n=1,m=2$, a configuration consists of one point
$p_1\in\HH$ and two ordered real points $q_1<q_2$. Quotienting by
$z\mapsto az+b$ with $a>0$ lets one normalize
\[
q_1=0,\qquad q_2=1.
\]
The remaining point $p_1$ varies in the upper half-plane. The graph
weight integrates the angle form over the compactified version of this
space.
\end{example}

\begin{theorem}[Kontsevich's explicit formality on $\R^{d}$
{\cite[\S6.4]{Kontsevich2003}}]\label{thm:explicit}
Restricted to bivectors, the maps
\[
U_{n}(\gamma_{1}, \dots, \gamma_{n})
= \sum_{\Gamma \in G_{n,2}}
w_{\Gamma}\, B_{\Gamma}(\gamma_{1}, \dots, \gamma_{n})
\qquad (\gamma_{k} \in \Gamma(\Lambda^{2}T_{\R^{d}}))
\]
extend $U_{\mathrm{HKR}}$ and, together with their analogues on graphs
with aerial vertices of higher outdegree (Remark~\ref{rem:bivectors-only}),
assemble into an $L_\infty$-quasi-isomorphism
$\Tpoly(\R^{d}) \to \Dpoly(\R^{d})$.
The $L_\infty$ identities follow from Stokes' theorem: for each
graph~$\Gamma$ appearing in a given identity, the closed form
$\bigwedge_{e} d\phi_{e}$ on the compactified configuration space
$\bar C_{n,m+1}^{+}$ has an integral that decomposes, via Stokes, into
contributions from the codimension-$1$ boundary strata. These strata
correspond bijectively to compositions of two admissible graphs, and
the resulting sum is precisely the $L_\infty$ relation.
\end{theorem}

\begin{example}[Constant Poisson tensors]
\label{ex:constant-poisson-graphs}
If the coefficients $\pi^{ij}$ are constant, then every graph with an
edge landing at an aerial vertex gives zero, because it differentiates
a coefficient of $\pi$. Only graphs whose edges land on ground vertices
survive. In this case Kontsevich's formula reduces to the Moyal product
of Example~\ref{ex:moyal-background}.
\end{example}

\begin{example}[Boundary strata and Stokes' theorem]
\label{ex:boundary-stokes}
When points in a configuration collide, the compactification
$\bar C_{n,m}$ records the limiting direction of the collision. These
codimension-$1$ boundary faces correspond to compositions of graphs.
Stokes' theorem says that the integral of an exact form over the
compactified space is zero. The resulting sum of boundary contributions
is precisely the $L_\infty$ identity for the maps $U_n$.
\end{example}

\section{Formalization path in Lean / Mathlib}\label{sec:lean}

This section is the working plan for Track~4. The aim is to make
Theorem~\ref{thm:formality} a Lean statement and then attack
Theorem~\ref{thm:explicit} layer by layer.

\begingroup
\emergencystretch=3em
\hbadness=10000

\subsection{Available infrastructure in Mathlib}\label{ssec:available}

Audited against Mathlib at commit \path{360da6fa66} (2026-06-18,
toolchain \path{leanprover/lean4:v4.32.0-rc1}). The following are in
Mathlib at a usable level.

\emph{Linear and graded algebra.}
\begin{itemize}[leftmargin=2em]
  \item Tensor and exterior algebras of a module, in
        \path{Mathlib.LinearAlgebra.TensorAlgebra.Basic}
        (with grading in \path{TensorAlgebra.Grading}) and
        \path{Mathlib.LinearAlgebra.ExteriorAlgebra.Basic}
        (grading, basis, and construction from alternating maps in
        the sibling files).
  \item Exterior powers $\Lambda^{k} V$ and alternating multilinear
        maps as first-class objects, in
        \path{Mathlib.LinearAlgebra.ExteriorPower} and
        \path{Mathlib.LinearAlgebra.Alternating}.
  \item Graded monoids, graded rings, graded algebras, and graded
        modules, in \path{Mathlib.Algebra.GradedMonoid},
        \path{Mathlib.Algebra.GradedMulAction},
        \path{Mathlib.Algebra.Module.GradedModule},
        \path{Mathlib.Algebra.DirectSum.Algebra},
        and \path{Mathlib.Algebra.DirectSum.Ring}.
  \item Lie rings, Lie algebras, and Lie modules in
        \path{Mathlib.Algebra.Lie.Basic}; graded Lie algebras in
        \path{Mathlib.Algebra.Lie.Graded}.
  \item Derivations of a commutative ring / algebra in
        \path{Mathlib.RingTheory.Derivation.Basic}
        (with the Lie-bracket instance in
        \path{Mathlib.RingTheory.Derivation.Lie}).
\end{itemize}

\emph{Formal power series.}
\begin{itemize}[leftmargin=2em]
  \item Univariate formal power series (used for the coefficient ring
        $\R[[\hbar]]$) in \path{Mathlib.RingTheory.PowerSeries.Basic},
        with substitution, order, inverses, and truncation in the
        sibling files.
  \item Multivariate power series in
        \path{Mathlib.RingTheory.MvPowerSeries.Basic}
        with substitution and evaluation in the sibling files.
\end{itemize}

\emph{Homological algebra.}
\begin{itemize}[leftmargin=2em]
  \item Chain and cochain complexes over an abelian category
        (\path{Mathlib.Algebra.Homology.HomologicalComplex},
        \path{Mathlib.Algebra.Homology.CochainComplexPlus}), together
        with quasi-isomorphisms
        (\path{Mathlib.Algebra.Homology.QuasiIso}) and homology
        functors. This is the dg framework we will wrap dg Lie
        algebras around.
\end{itemize}

\emph{Smooth manifolds, tangent bundles, vector fields.}
\begin{itemize}[leftmargin=2em]
  \item Smooth manifolds via
        \path{Mathlib.Geometry.Manifold.IsManifold.Basic}, built on
        the \emph{ModelWithCorners} abstraction. So manifolds with
        corners are supported at the level of models. Interior and
        boundary of such a manifold, and the fact that they partition
        it, live in
        \path{Mathlib.Geometry.Manifold.IsManifold.InteriorBoundary}.
  \item Tangent bundle as a $C^{n}$ vector bundle
        (\path{Mathlib.Geometry.Manifold.VectorBundle.Tangent})
        and smooth sections of a vector bundle
        (\path{Mathlib.Geometry.Manifold.VectorBundle.ContMDiffSection}).
  \item Point derivations of $\Cinfty(M)$ and the differential of a
        smooth map, in
        \path{Mathlib.Geometry.Manifold.DerivationBundle}.
  \item Lie bracket of vector fields on a manifold, in
        \path{Mathlib.Geometry.Manifold.VectorField.LieBracket}.
\end{itemize}

\emph{Differential forms and calculus.}
\begin{itemize}[leftmargin=2em]
  \item Exterior derivative of a differential form on a normed space
        in \path{Mathlib.Analysis.Calculus.DifferentialForm.Basic}
        (Kudryashov--Lindauer, 2025), with the Cartan-style pairing
        against vector fields in \path{DifferentialForm.VectorField}.
        Only defined on normed spaces so far: the file explicitly
        lists ``bundled $C^{r}$-smooth $n$-forms on manifolds'' as a
        TODO.
\end{itemize}

\subsection{Missing infrastructure}\label{ssec:missing}

The audit above shows the following are absent from Mathlib. Each
gap becomes a workstream in Section~\ref{sec:milestones}.

\emph{Homotopical algebra.}
\begin{itemize}[leftmargin=2em]
  \item $L_\infty$-algebras and $L_\infty$-morphisms with Koszul
        signs. No \path{LInftyAlgebra}, no \path{MaurerCartan},
        nothing in Mathlib on either.
  \item Maurer--Cartan elements of a dg Lie or $L_\infty$-algebra,
        and gauge equivalence over $\hbar \K[[\hbar]]$.
  \item Operads and algebras over operads (needed only for the
        Tamarkin alternative in Section~\ref{ssec:alternatives}).
\end{itemize}

\emph{Hochschild theory of associative algebras.}
\begin{itemize}[leftmargin=2em]
  \item Hochschild cochain complex of an associative algebra $A$
        (multilinear maps $A^{\otimes n} \to A$ with the Hochschild
        differential). A repository-wide search returns exactly one
        mention of ``Hochschild'', a comment on the Hochschild--Serre
        spectral sequence in group cohomology; there is no Hochschild
        homology or cohomology of associative algebras in Mathlib.
  \item Gerstenhaber bracket on that complex; the associated dg Lie
        structure; the fact that MC elements are deformations of the
        product.
\end{itemize}

\emph{Poisson / polyvector-field side.}
\begin{itemize}[leftmargin=2em]
  \item Polyvector fields $\Gamma(\Lambda^{n+1} T_M)$ as a graded
        $\Cinfty(M)$-module. The pieces exist (exterior power of a
        module; tangent bundle; smooth sections) but not the graded
        object $\Tpoly(M)$.
  \item Schouten--Nijenhuis bracket on $\Tpoly(M)$, its Jacobi
        identity, and the vanishing of $\Sch{\cdot}{\cdot}$ on
        functions.
  \item Poisson manifolds, Poisson brackets, the equivalence
        between Poisson bivectors and Lie brackets on $\Cinfty(M)$.
        Grep for ``Poisson'' in \path{Mathlib/Geometry} returns
        nothing; the existing hits under \path{Analysis} are the
        Poisson kernel and Poisson summation, unrelated.
\end{itemize}

\emph{Polydifferential operators on $\Cinfty(M)$.}
\begin{itemize}[leftmargin=2em]
  \item Polydifferential operators as a concrete subtype of
        $\Cinfty(M)^{\otimes\bullet} \to \Cinfty(M)$. On $\R^{d}$
        with polynomial coefficients this is a routine algebraic
        subtype; on a general manifold it requires jet bundles.
  \item Star products on $(M,\pi)$ and gauge equivalence between
        them.
\end{itemize}

\emph{Forms, integration, corners.}
\begin{itemize}[leftmargin=2em]
  \item Differential forms on manifolds (Mathlib has them only on
        normed spaces, per the file \path{DifferentialForm/Basic.lean}
        TODO).
  \item Integration of differential forms on an oriented manifold,
        including the case of a manifold with corners.
  \item Stokes' theorem for manifolds with corners. The abstract
        model \emph{ModelWithCorners} is available and boundary points
        are defined, but there is no integration theory that uses
        them.
\end{itemize}

\emph{Configuration spaces.}
\begin{itemize}[leftmargin=2em]
  \item Kontsevich admissible graphs $G_{n,m}$ as a finite decidable
        type.
  \item Fulton--MacPherson--Axelrod--Singer compactifications
        $\bar C_{n,m}$, their manifold-with-corners structure, and
        the identification of codimension-$1$ boundary strata with
        compositions of smaller configuration spaces.
  \item The Kontsevich angle 1-form $d\phi_{e}$ on
        $\bar C_{n,m}^{+}$ and the weight integrals $w_{\Gamma}$.
\end{itemize}

\endgroup

\begingroup
\emergencystretch=3em
\hbadness=10000

\subsection{Task decomposition for the initial team}\label{ssec:tasks}

The six workstreams of Section~\ref{ssec:missing} split into atomic
tasks. Each task fixes a Lean module path, a mathematical statement,
its dependencies among other tasks, and a target sorry count. The
sorry-count target is $0$ for algebraic tasks and equals the number
of documented axiom blocks for tasks that isolate the analytic input.
Effort estimates are $S$ (one to two weeks), $M$ (two to four weeks),
$L$ (four to eight weeks) for a Lean-fluent contributor. The task
graph is laid out so that Workstream~A, and the first task of each
other workstream, can be started in parallel on day one and produce a
visible artefact within two weeks.

\paragraph{Workstream A. Combinatorial and polynomial scaffolding.}
Foundational Lean-typeclass work. No non-trivial mathematics.
Serves as the entry point for contributors calibrating their Lean.
\begin{itemize}[leftmargin=2em]
  \item[\textbf{A1}]
    \defi{Admissible graphs as a decidable finite type.} $S$.
    Module \path{Kontsevich.Combinatorics.AdmissibleGraph}. Implement
    the structure of Definition~\ref{def:graph} with
    \texttt{DecidableEq} and a \texttt{Fintype} instance on
    $G_{n,m}$. Deliverable: \texttt{\#eval} enumerates $G_{1,2}$
    (one graph), $G_{2,2}$ (four graphs after edge labelling), and
    $G_{3,2}$. Depends on: none.
  \item[\textbf{A2}]
    \defi{Polynomial ring on $\R^{d}$ as an ambient algebra.} $S$.
    Module \path{Kontsevich.Polynomial.RdPoly}. Wrap
    \path{Mathlib.Algebra.MvPolynomial.Basic} with the smooth-function
    interpretation: coordinate derivations
    $\partial_{i}$ as \path{MvPolynomial.derivation}, and a lemma
    that they commute pairwise. Depends on: none.
  \item[\textbf{A3}]
    \defi{Polynomial polyvector fields on $\R^{d}$.} $S$.
    Module \path{Kontsevich.Polynomial.Tpoly}. Define
    $\Tpoly^{n}(\R^{d})_{\mathrm{poly}}$ as
    $\R[x_{1},\dots,x_{d}] \otimes \Lambda^{n+1} \R^{d}$ using
    \path{Mathlib.LinearAlgebra.ExteriorPower}. Depends on: A2.
  \item[\textbf{A4}]
    \defi{Polynomial polydifferential operators on $\R^{d}$.} $M$.
    Module \path{Kontsevich.Polynomial.Dpoly}. Define
    $\Dpoly^{n}(\R^{d})_{\mathrm{poly}}$ as multilinear maps
    $\R[x]^{\otimes(n+1)} \to \R[x]$ whose action factors as a
    finite polynomial combination of $\partial^{\alpha}$'s. Depends
    on: A2.
  \item[\textbf{A5}]
    \defi{The operator $B_{\Gamma}$ in polynomial coefficients.}
    $M$. Module \path{Kontsevich.Polynomial.BGamma}. Implement
    Definition~\ref{def:BGamma} for
    $\gamma_{k} \in \Tpoly^{1}(\R^{d})_{\mathrm{poly}}$ and
    $f_{j} \in \R[x]$. Depends on: A1, A3, A4.
  \item[\textbf{A6}]
    \defi{Kontsevich weights as an axiomatised black box.} $S$.
    Module \path{Kontsevich.Analysis.WeightAxiom}. Declare
    $w \colon G_{n,2} \to \R$ as an \texttt{axiom} for each $n$, with
    the values for $n = 1, 2$ substituted by the closed-form rationals
    from \cite[\S8.3]{Kontsevich2003}. Weights for $n \ge 3$ stay
    abstract. Depends on: A1.
\end{itemize}

\paragraph{Workstream B. Graded and Lie infrastructure.}
Bridges Mathlib's existing typeclasses to the dg Lie structures the
project needs. Runs in parallel with A.
\begin{itemize}[leftmargin=2em]
  \item[\textbf{B1}]
    \defi{dg Lie algebra wrapper on Mathlib.} $S$ to $M$. Module
    \path{Kontsevich.Algebra.DGLie}. Assemble a \texttt{DGLA}
    structure from \path{Mathlib.Algebra.Homology.HomologicalComplex}
    over the category of vector spaces and
    \path{Mathlib.Algebra.Lie.Basic}. Prove Jacobi and the Leibniz
    rule for $d$ as instance obligations. Depends on: none.
  \item[\textbf{B2}]
    \defi{Schouten--Nijenhuis bracket on
          $\Tpoly(\R^{d})_{\mathrm{poly}}$.} $M$. Module
    \path{Kontsevich.Polynomial.Schouten}. Implement
    $\Sch{\cdot}{\cdot}$ via the local-coordinate formula, prove
    graded antisymmetry and the graded Jacobi identity. Deliverable:
    theorems \texttt{schouten\_antisymm} and \texttt{schouten\_jacobi}.
    Depends on: A2, A3.
  \item[\textbf{B3}]
    \defi{Gerstenhaber bracket on
          $\Dpoly(\R^{d})_{\mathrm{poly}}$.} $M$. Module
    \path{Kontsevich.Polynomial.Gerstenhaber}. Implement
    $\Ger{\cdot}{\cdot}$ via the sum-over-insertions formula
    (\cite{Gerstenhaber1963}, \S6), prove graded antisymmetry and
    graded Jacobi. Depends on: A2, A4.
  \item[\textbf{B4}]
    \defi{Hochschild differential and $d_{\mathrm{H}}^{2}=0$.} $S$.
    Module \path{Kontsevich.Polynomial.Hochschild}. Set
    $d_{\mathrm{H}} := \Ger{\mu}{\cdot}$ where $\mu$ is the pointwise
    product. Prove $d_{\mathrm{H}}^{2} = 0$ by expanding
    $\Ger{\Ger{\mu}{\mu}}{\cdot}$. Depends on: B3.
  \item[\textbf{B5}]
    \defi{Two dg Lie instances.} $S$. Module
    \path{Kontsevich.Polynomial.DGLieInstances}. Register
    $\Tpoly(\R^{d})_{\mathrm{poly}}$ (with zero differential) and
    $\Dpoly(\R^{d})_{\mathrm{poly}}$ (with differential
    $d_{\mathrm{H}}$) as \texttt{DGLA} instances via B1. Depends on:
    B1, B2, B4.
\end{itemize}

\paragraph{Workstream C. Hochschild--Kostant--Rosenberg.}
Bridges the two dg Lie algebras of Workstream~B. Contains the first
non-trivial theorem the team can point at.
\begin{itemize}[leftmargin=2em]
  \item[\textbf{C1}]
    \defi{$U_{\mathrm{HKR}}$ as a graded linear map.} $S$. Module
    \path{Kontsevich.Polynomial.HKR}. Define
    $U_{\mathrm{HKR}} \colon \Tpoly(\R^{d})_{\mathrm{poly}}
    \to \Dpoly(\R^{d})_{\mathrm{poly}}$ by
    Theorem~\ref{thm:HKR}. Depends on: A3, A4.
  \item[\textbf{C2}]
    \defi{$U_{\mathrm{HKR}}$ is a chain map.} $M$. Same module.
    Prove $d_{\mathrm{H}} \circ U_{\mathrm{HKR}} = 0$ (the domain has
    zero differential). Depends on: B4, C1.
  \item[\textbf{C3}]
    \defi{$U_{\mathrm{HKR}}$ is a quasi-isomorphism on
          $\R^{d}$-polynomials.} $M$ to $L$. Module
    \path{Kontsevich.Polynomial.HKRQuasiIso}. Compute Hochschild
    cohomology of $\R[x_{1},\dots,x_{d}]$ using the Koszul
    resolution and match it to $\Tpoly(\R^{d})_{\mathrm{poly}}$.
    Depends on: C2.
\end{itemize}

\paragraph{Workstream D. $L_\infty$ machinery and Maurer--Cartan.}
The sign-heavy homotopical algebra core. Sign conventions are fixed
in D1 and cited by label everywhere else in the project.
\begin{itemize}[leftmargin=2em]
  \item[\textbf{D1}]
    \defi{$L_\infty$-algebra typeclass with Koszul signs.} $L$.
    Module \path{Kontsevich.LInfty.Algebra}. Implement
    Definition~\ref{def:Linfty}. Fix the sign convention of
    \cite[\S3.4]{Kontsevich2003} in an \texttt{assumption} block and
    prove the $n=1,2,3$ special cases. Depends on: none.
  \item[\textbf{D2}]
    \defi{$L_\infty$-morphism typeclass.} $M$. Module
    \path{Kontsevich.LInfty.Morphism}. Implement
    Definition~\ref{def:Linfty-morphism} and the
    quasi-isomorphism predicate on $U_{1}$. Depends on: D1.
  \item[\textbf{D3}]
    \defi{Embedding of dg Lie algebras into $L_\infty$-algebras.} $S$.
    Module \path{Kontsevich.LInfty.FromDGLA}. Prove that a
    \texttt{DGLA} is an $L_\infty$-algebra with $Q_{n} = 0$ for
    $n \ge 3$. Depends on: B1, D1.
  \item[\textbf{D4}]
    \defi{Maurer--Cartan elements over $\hbar \K[[\hbar]]$.} $M$.
    Module \path{Kontsevich.LInfty.MaurerCartan}. Implement
    Definition~\ref{def:MC}. Prove that MC elements of a dg Lie
    algebra satisfy $d\alpha + \tfrac{1}{2}[\alpha,\alpha] = 0$.
    Depends on: D1.
  \item[\textbf{D5}]
    \defi{Gauge equivalence of MC elements.} $M$ to $L$. Module
    \path{Kontsevich.LInfty.Gauge}. Implement the exponential action
    of degree-$0$ elements and prove that gauge equivalence is an
    equivalence relation on MC elements. Depends on: D4.
  \item[\textbf{D6}]
    \defi{MC transport under $L_\infty$-morphisms.} $M$. Module
    \path{Kontsevich.LInfty.MCTransport}. Prove that
    $\alpha \mapsto \sum_{n\ge 1} \tfrac{1}{n!} U_{n}(\alpha^{\wedge n})$
    sends MC to MC. Depends on: D2, D4.
\end{itemize}

\paragraph{Workstream E. Kontsevich formula on $\R^{d}$.}
Assembles A--D into the polynomial-coefficient formality statement.
\begin{itemize}[leftmargin=2em]
  \item[\textbf{E1}]
    \defi{The map $U_{n}$ in polynomial coefficients.} $M$. Module
    \path{Kontsevich.Formula.U}. Define
    $U_{n} := \sum_{\Gamma \in G_{n,2}} w_{\Gamma}\, B_{\Gamma}$
    as a map
    $\bigl(\Tpoly^{1}(\R^{d})_{\mathrm{poly}}\bigr)^{\wedge n}
    \to \Dpoly^{1}(\R^{d})_{\mathrm{poly}}$. Prove $U_{1}$ equals
    $U_{\mathrm{HKR}}$ restricted to bivectors. Depends on: A5, A6,
    C1.
  \item[\textbf{E2}]
    \defi{The Moyal product from the Kontsevich formula.} $M$.
    Module \path{Kontsevich.Formula.Moyal}. Specialise E1 to
    constant $\pi$ and prove the resulting star product equals the
    Moyal product of Example~\ref{ex:moyal-background}. Prove
    associativity directly by graph enumeration for $n \le 3$ and
    inductively for general $n$. Depends on: A5, E1.
  \item[\textbf{E3}]
    \defi{$L_\infty$ identities modulo boundary axioms.} $L$.
    Module \path{Kontsevich.Formula.LInftyIds}. Prove that the maps
    $U_{n}$ of E1 satisfy the $L_\infty$ identities of D1 modulo an
    explicit \texttt{axiom} block encoding the boundary identities
    of Theorem~\ref{thm:explicit}. Sorry count: $0$; axiom count:
    one per required boundary vanishing. Depends on: B5, D3, D6, E1.
  \item[\textbf{E4}]
    \defi{Star-product corollary.} $M$. Module
    \path{Kontsevich.Formula.StarProduct}. Deduce
    Corollary~\ref{cor:quantization} for polynomial Poisson tensors
    on $\R^{d}$: the series
    $f\star g = fg + \sum_{n\ge 1} \tfrac{\hbar^{n}}{n!}
    U_{n}(\pi,\dots,\pi)(f,g)$ is associative. Depends on: D6, E3.
\end{itemize}

\paragraph{Workstream F. Analytic input (phase two).}
Discharges the axioms of A6 and E3. This workstream is deferred
until Workstreams A--E deliver Milestones~1--8 of
Section~\ref{sec:milestones}; documented here so that the axiom
count of E3 is bounded by a concrete list.
\begin{itemize}[leftmargin=2em]
  \item[\textbf{F1}]
    \defi{Configuration space $C_{n,m}$ as a smooth manifold.} $L$.
    Module \path{Kontsevich.Analytic.Config}. Define $C_{n,m}$ and
    equip it with the smooth structure from
    Definition~\ref{def:weight}. Depends on: A1.
  \item[\textbf{F2}]
    \defi{The Kontsevich angle 1-form $d\phi_{e}$.} $M$. Module
    \path{Kontsevich.Analytic.AngleForm}. Depends on: F1.
  \item[\textbf{F3}]
    \defi{Compactification $\bar C_{n,m}$ as a manifold with
          corners.} $L$. Module \path{Kontsevich.Analytic.Compactify}.
    Build the Fulton--MacPherson--Axelrod--Singer compactification.
    Identify the codimension-$1$ boundary strata. Depends on: F1.
  \item[\textbf{F4}]
    \defi{Stokes' theorem on $\bar C_{n,m}^{+}$ and boundary
          vanishing.} $L$. Module \path{Kontsevich.Analytic.Stokes}.
    Discharge the axiom block of E3. Depends on: F2, F3.
  \item[\textbf{F5}]
    \defi{Weight computation for low $n$.} $S$ to $M$. Module
    \path{Kontsevich.Analytic.WeightsLow}. Replace the axiomatised
    values of A6 for $n = 1, 2$ with theorems that evaluate the
    integral $w_{\Gamma}$. Depends on: F2, F3.
\end{itemize}

\paragraph{Suggested allocation.}
A team of six to eight contributors of mixed Lean fluency covers the
six workstreams. A workable initial mapping:
two contributors on Workstream~A (foundational Lean, junior
friendly), one on B1 shared with two more on B2--B5 (mid-level Lean
plus one algebraist per bracket), one on Workstream~C (algebraic
depth), one to two on Workstream~D (sign-heavy, senior Lean), and
one shared coordinator on Workstream~E. Workstream~F is deferred.
The task graph has three cut points that gate visible team
progress: A5+A6 (a runnable $B_{\Gamma}$ on graphs), B2+B3
(the two brackets), and E2 (the Moyal product as a Kontsevich
computation). Each cut point yields a self-contained artefact that
can be reviewed independently.

\endgroup

\subsection{Layered plan}\label{ssec:plan}

\paragraph{Layer 0 (Combinatorial scaffolding).}
Implement \texttt{AdmissibleGraph (n m : Nat)} from
Definition~\ref{def:graph} as a finite, decidable structure. Build the
polynomial-coefficient version of $B_{\Gamma}$ on $\R^{d}$ acting on
polynomial functions. Axiomatize $w_{\Gamma} : \mathbb{Q}_{\Gamma} \to \R$
as a real-valued constant per graph, postponing its integral definition.
This layer is purely finite and decidable. It is testable: for small
$n, m$ one can check the Kontsevich formula returns the Moyal product
on the constant Poisson structure.

\paragraph{Layer 1. Algebraic $L_\infty$.}
Define \texttt{LInftyAlgebra (V : Type*) [GradedModule V]} with the
maps $Q_{n}$ of Definition~\ref{def:Linfty} and the generalized Jacobi
identities. Define \texttt{LInftyMorphism} per
Definition~\ref{def:Linfty-morphism}. Construct the dg-Lie embedding
$\mathrm{DGLA} \hookrightarrow L_\infty\mathrm{-Alg}$. Define Maurer--Cartan
elements and gauge equivalence over $\hbar \K[[\hbar]]$.
Sign conventions follow \cite[\S4]{Kontsevich2003}; they must be
declared once and reused.

\paragraph{Layer 2. Concrete dg-Lie algebras.}
For $V = \R^{d}$, define $\Tpoly(\R^{d})_{\mathrm{poly}}$ as polynomial
polyvector fields with the Schouten--Nijenhuis bracket; this avoids
analysis. Define $\Dpoly(\R^{d})_{\mathrm{poly}}$ as polydifferential
operators with polynomial coefficients, with the Gerstenhaber bracket
and Hochschild differential. Prove $U_{\mathrm{HKR}}$ is a chain map
(Theorem~\ref{thm:HKR}) in this polynomial setting.

\paragraph{Layer 3. Formality.}
With the axiomatized weights of Layer~0 and the algebraic infrastructure
of Layers~1--2, prove that the maps $U_{n}$ of Theorem~\ref{thm:explicit}
define an $L_\infty$-morphism modulo a single hypothesis: the boundary
integrals on $\bar C_{n,m+1}^{+}$ vanish in the combinations required
by the $L_\infty$ identities. The hypothesis is the only analytic
input; everything else is finite combinatorics with rational
coefficients.

\paragraph{Layer 4. Discharging the analytic hypothesis.}
Two routes are available.
\begin{enumerate}[label=(\alph*),leftmargin=2em]
  \item Build the integration theory: compactified configuration
        spaces, integration of forms on manifolds with corners,
        Stokes' theorem with corners. Then prove the boundary
        vanishing directly.
  \item Replace the analytic input by an axiom block citing
        \cite[\S6]{Kontsevich2003} and document the gap. This is the
        recommended starting point.
\end{enumerate}

\subsection{Subproblems that can be closed without analysis}
\label{ssec:no-analysis}

Two reduced problems are fully algebraic and can be formalized today:
\begin{itemize}[leftmargin=2em]
  \item \defi{Linear Poisson structures on $\mathfrak{g}^{*}$.} The
        restriction of $U$ to linear Poisson tensors on the dual of a
        finite-dimensional Lie algebra $\mathfrak{g}$ equals the
        symmetrization map $\mathrm{Sym}(\mathfrak{g}) \to
        \mathcal{U}(\mathfrak{g})$ composed with the Duflo isomorphism
        (Kathotia \cite[Thm.~1.1]{Kathotia2000}, Shoikhet
        \cite{Shoikhet2001}). The PBW theorem and the Duflo factor
        can be formalized without configuration-space integrals.
  \item \defi{The Moyal product.} For constant Poisson structures on
        $\R^{d}$, all weights $w_{\Gamma}$ with $\Gamma$ having edges
        among aerial vertices vanish, and the formula reduces to the
        exponential of $\hbar \pi^{ij} \partial_{i} \otimes
        \partial_{j} / 2$. This is the natural first end-to-end
        verification target.
\end{itemize}

\subsection{Challenges}\label{ssec:challenges}

\begin{itemize}[leftmargin=2em]
  \item \defi{Sign conventions.} Koszul signs in the $L_\infty$
        identities are the dominant source of errors in any
        formalization of homotopical algebra. They must be fixed in
        Layer~1 and never revisited.
  \item \defi{Transcendental weights.} The numbers $w_{\Gamma}$ are
        periods of differential forms on manifolds with corners. Some
        are known closed form (e.g.\ wheel weights are computed in
        \cite{Shoikhet2001}); most are not. Axiomatization is required
        unless one ports the integration theory.
  \item \defi{Polydifferential operators.} A faithful Lean type for
        polydifferential operators on $\Cinfty(M)$ requires a careful
        treatment of jet bundles or local coordinate charts. Starting
        with polynomial coefficients on $\R^{d}$ side-steps this.
  \item \defi{Gauge equivalence.} The classification statement in
        Corollary~\ref{cor:quantization} requires gauge equivalence of
        MC elements modulo $\hbar$, which requires the full $L_\infty$
        machinery, including the homotopy transfer needed for
        equivalences.
  \item \defi{Globalization.} Kontsevich's $\R^{d}$ formula globalizes
        to manifolds by Fedosov-style methods
        (Cattaneo--Felder--Tomassini \cite{CFT}, Dolgushev
        \cite{Dolgushev2005}). This is a second-stage project. It
        depends on a connection / formal geometry layer in Mathlib
        that does not yet exist.
\end{itemize}

\subsection{Alternative approaches}\label{ssec:alternatives}

\begin{itemize}[leftmargin=2em]
  \item \defi{Tamarkin's proof} \cite{Tamarkin}. Replaces the
        configuration-space integrals by the formality of the little
        disks operad $E_{2}$ and a choice of Drinfeld associator. This
        is more algebraic and avoids analysis, but requires operad
        theory in Mathlib, which is also missing.
  \item \defi{Cattaneo--Felder Poisson sigma model} \cite{CF}. Gives
        $w_{\Gamma}$ as expectation values of a topological field
        theory on the disk. Not directly formalizable, but a useful
        guide for interpreting the boundary identities of
        Theorem~\ref{thm:explicit}.
  \item \defi{Manchon--Torossian} \cite{MT}. Refines the tangent
        cohomology of the Kontsevich star product and gives sharper
        statements at the level of cup products. A target for
        formalization once the base layer is in place.
\end{itemize}

\section{Concrete first milestones}\label{sec:milestones}

The proposed sequence of milestones, each independently verifiable:
\begin{enumerate}[leftmargin=2em]
  \item \texttt{AdmissibleGraph} type with decidable equality and
        enumeration; unit tests against \cite[\S6.1]{Kontsevich2003}.
  \item \texttt{LInftyAlgebra} structure with Koszul signs; the
        embedding $\mathrm{DGLA} \hookrightarrow L_\infty\mathrm{-Alg}$.
  \item Schouten--Nijenhuis bracket on polynomial polyvector fields on
        $\R^{d}$; Jacobi identity proved.
  \item Gerstenhaber bracket on polynomial polydifferential operators
        on $\R^{d}$; Hochschild differential proved square-zero.
  \item Theorem~\ref{thm:HKR} as a Lean theorem in the polynomial
        setting.
  \item Moyal product as the Kontsevich formula evaluated on constant
        $\pi$; associativity proved by graph enumeration.
  \item PBW-symmetrization formula for linear $\pi$ on
        $\mathfrak{g}^{*}$; Duflo factor handled as an axiom on first
        pass.
  \item Statement of Theorem~\ref{thm:formality} with axiomatized
        weights; the $L_\infty$ identities reduced to the boundary
        hypothesis.
\end{enumerate}

Milestones~1--6 are fully algebraic and have no analytic axioms.
Milestone~7 has one axiom (the Duflo factor) that is independently
documentable. Milestone~8 is the cleanest target for a first paper.

\begingroup
\emergencystretch=3em
\hbadness=10000
\begin{thebibliography}{99}

\bibitem{BFFLS}
F.~Bayen, M.~Flato, C.~Fronsdal, A.~Lichnerowicz, D.~Sternheimer,
\emph{Deformation theory and quantization. I. Deformations of
symplectic structures}, Ann. Physics \textbf{111} (1978), 61--110.

\bibitem{CF}
A.~S.~Cattaneo, G.~Felder,
\emph{A path integral approach to the Kontsevich quantization formula},
Comm. Math. Phys. \textbf{212} (2000), 591--611.

\bibitem{CFT}
A.~S.~Cattaneo, G.~Felder, L.~Tomassini,
\emph{From local to global deformation quantization of Poisson
manifolds}, Duke Math. J. \textbf{115} (2002), 329--352.

\bibitem{Dolgushev2005}
V.~Dolgushev,
\emph{Covariant and equivariant formality theorems},
Adv. Math. \textbf{191} (2005), 147--177.

\bibitem{Gerstenhaber1963}
M.~Gerstenhaber,
\emph{The cohomology structure of an associative ring},
Ann. of Math. (2) \textbf{78} (1963), 267--288.

\bibitem{Kathotia2000}
V.~Kathotia,
\emph{Kontsevich's universal formula for deformation quantization and
the Campbell--Baker--Hausdorff formula},
Internat. J. Math. \textbf{11} (2000), 523--551.

\bibitem{Kontsevich1999}
M.~Kontsevich,
\emph{Operads and motives in deformation quantization},
Lett. Math. Phys. \textbf{48} (1999), 35--72.

\bibitem{Kontsevich2003}
M.~Kontsevich,
\emph{Deformation quantization of Poisson manifolds},
Lett. Math. Phys. \textbf{66} (2003), 157--216.
arXiv:q-alg/9709040.

\bibitem{MT}
D.~Manchon, C.~Torossian,
\emph{Cohomologie tangente et cup-produit pour la quantification de
Kontsevich}, Ann. Math. Blaise Pascal \textbf{10} (2003), 75--106.

\bibitem{Shoikhet2001}
B.~Shoikhet,
\emph{Vanishing of the Kontsevich integrals of the wheels},
Lett. Math. Phys. \textbf{56} (2001), 141--149.

\bibitem{Tamarkin}
D.~Tamarkin,
\emph{Another proof of M.~Kontsevich formality theorem},
arXiv:math/9803025 (1998).

\end{thebibliography}
\endgroup

\end{document}
